\chapter{颜色与光源}\label{ch:g11:color-light-sources}

问一位物理学家番茄是什么颜色，得到的却是一个反问：
在哪种光下？番茄只漫射它所接收的光，灯
只发出其光谱所含的光，而眼睛把到达的任何光
都压缩为三个数。颜色是光源、
物体与眼睛之间的协商——本章将逐一审视这三方。

\section{以三个数看颜色}

\begin{definition}[视锥细胞与三色视觉]\label{def:g11:color-light-sources:cones}
视网膜（见\cref{ch:g11:lenses-and-eye}）含有两类感光细胞：
\emph{视杆细胞}\index{视杆细胞}，对微弱的光有反应，却不分辨波长
（夜视时呈灰色），另一类是
\emph{视锥细胞}\index{视锥细胞}，共三种，分别在
$\qty{420}{nm}$（蓝）、$\qty{530}{nm}$（绿）和 $\qty{560}{nm}$
（红）附近最敏感。人的颜色视觉属于
\emph{三色的}\index{三色视觉}：无论光谱如何，
大脑接收到的只有三种视锥细胞的响应，颜色
便是大脑由这一三元组所形成的感受。
\end{definition}

\begin{definition}[颜色即感知]\label{def:g11:color-light-sources:perception}
因此，颜色是一种\emph{感知}，而非仅由光本身决定的性质：
光谱给出每个波长处的光强，而眼睛只保留
三个数。光谱不同而三种视锥细胞响应相同的光
看上去完全相同：它们即为\emph{同色异谱}\index{同色异谱}。
\end{definition}

\begin{example}[两种黄色]\label{ex:g11:color-light-sources:yellows}
钠灯的 $\qty{589}{nm}$ 谱线几乎等量地激发红、绿
两类视锥细胞，大脑便判断为\emph{黄色}。一块屏幕点亮
其红色（$\qty{610}{nm}$）和绿色（$\qty{540}{nm}$）子像素，
所呈现的黄色与之无异，而 $\qty{589}{nm}$ 附近却无光。棱镜
能立即将两者分开，而眼睛永远做不到。
\end{example}

\section{光的相加}

\begin{definition}[加色混合]\label{def:g11:color-light-sources:additive}
将光束叠加于白色屏幕上，即是将它们的光谱相加：此即
\emph{加色混合}\index{加色混合}。选取恰当的三种光，
即红、绿、蓝——这三种\emph{三原色}\index{三原色}，
几乎可混合出所有可感知的颜色：
R $+$ G $=$ 黄色，R $+$ B $=$ 品红，G $+$ B $=$ 青色，
以及 R $+$ G $+$ B $=$ 白色。
\end{definition}

\begin{proposition}[互补色]\label{prop:g11:color-light-sources:complementary}
若两束光之和被感知为白色，则称它们为
\emph{互补色}\index{互补色}。三对颜色
黄/蓝、青/红与品红/绿均互补。
\end{proposition}

\begin{proof}
黄色即 R $+$ G；加上蓝色，三原色便合成白色。另外
两对的论证与此相同，每次只换一种原色。
\end{proof}

\begin{example}[屏幕]\label{ex:g11:color-light-sources:screen}
手机屏幕的每个像素由三个子像素组成——红、绿、蓝，
宽度都小于 $\qty{0.1}{mm}$，近得眼睛无法分辨，故
它们的光在视网膜上相加：黄色为 R 与 G 同时点亮，橙色为 R 满功率、
G 半功率，白色则三者全亮。用放大镜观察
白色区域，便能看出其中的窍门：处处没有白色，只有 R、G、B。
\end{example}

\begin{omfigure}
\begin{tikzpicture}[scale=0.72, font=\small]
  \fill[red] (90:0.8) circle (1.25);
  \fill[green!70!black] (210:0.8) circle (1.25);
  \fill[blue] (330:0.8) circle (1.25);
  \begin{scope}\clip (90:0.8) circle (1.25);\fill[yellow] (210:0.8) circle (1.25);\end{scope}
  \begin{scope}\clip (90:0.8) circle (1.25);\fill[magenta] (330:0.8) circle (1.25);\end{scope}
  \begin{scope}\clip (210:0.8) circle (1.25);\fill[cyan] (330:0.8) circle (1.25);\end{scope}
  \begin{scope}\clip (90:0.8) circle (1.25);\clip (210:0.8) circle (1.25);
    \fill[white] (330:0.8) circle (1.25);\end{scope}
  \node at (90:2.5) {红};\node at (210:2.5) {绿};\node at (330:2.5) {蓝};
\end{tikzpicture}\qquad
\begin{tikzpicture}[scale=0.72, font=\small]
  \fill[cyan] (90:0.8) circle (1.25);
  \fill[magenta] (210:0.8) circle (1.25);
  \fill[yellow] (330:0.8) circle (1.25);
  \begin{scope}\clip (90:0.8) circle (1.25);\fill[blue] (210:0.8) circle (1.25);\end{scope}
  \begin{scope}\clip (90:0.8) circle (1.25);\fill[green!70!black] (330:0.8) circle (1.25);\end{scope}
  \begin{scope}\clip (210:0.8) circle (1.25);\fill[red] (330:0.8) circle (1.25);\end{scope}
  \begin{scope}\clip (90:0.8) circle (1.25);\clip (210:0.8) circle (1.25);
    \fill[black] (330:0.8) circle (1.25);\end{scope}
  \node at (90:2.5) {青};\node at (210:2.5) {品红};\node at (330:2.5) {黄};
\end{tikzpicture}

{\small 重叠光束的加色混合（左）与
白光下叠放滤光片的减色混合（右）。}
\end{omfigure}

\section{从白光中减去色光}

\begin{definition}[减色混合]\label{def:g11:color-light-sources:subtractive}
\emph{滤光片}\index{滤光片}（有色玻璃、凝胶、一层墨水或
颜料）透过其所接收光谱的一部分，并吸收其余部分。
通过从白光中\emph{去除}若干波段来产生颜色，称为
\emph{减色混合}\index{减色混合}。其
三原色即加色三原色的互补色，各自去除一个波段：
青色 $=$ 白色 $-$ R，品红 $=$ 白色 $-$ G，黄色 $=$ 白色 $-$ B。
叠放的滤光片共同相减：青色 $+$ 黄色得到绿色，青色
$+$ 品红得到蓝色，品红 $+$ 黄色得到红色，三者叠加为黑色。印刷
与绘画都以此方式实现，用的是 C、M、Y 墨水。
\end{definition}

\begin{method}[光谱跟踪法]\label{met:g11:color-light-sources:bookkeeping}
要预测经过任意光源、滤光片与物体的组合后所呈现的颜色，
应追踪三个原色波段：将光源的光写作它在 R、G、B 中的
成分；在每个滤光片或表面处删去被吸收的波段，
并说出留存之光的颜色（无剩余 $=$ 黑色）。
\end{method}

\begin{example}[滤光片串联]\label{ex:g11:color-light-sources:chain}
白光 $(R, G, B)$ 通过一块黄色滤光片：蓝色波段
被吸收，剩下 $(R, G)$——黄色。随后，青色滤光片吸收
红色波段，剩下 $(G)$：绿色。交换两块滤光片的次序
并不改变结果：每块滤光片只吸收自己的波段，与次序无关。
\end{example}

\section{物体的颜色}

\begin{definition}[漫射与物体颜色]\label{def:g11:color-light-sources:objectcolor}
不透明物体吸收所接收光的一部分，并
将其余部分\emph{漫射}\index{漫射}——即向各个方向再发射，
其中包括射向我们眼睛的方向。
\emph{物体的颜色}\index{物体的颜色}是物体所漫射光的颜色，
而该光须\emph{在给定照明下}产生：物体提供吸收特性，
光源提供原料。换一盏灯，颜色便随之改变。
\end{definition}

\begin{example}[红色衬衫]\label{ex:g11:color-light-sources:redshirt}
衬衫在日光下看起来是红色：它吸收绿色和蓝色波段，
并漫射红色波段。在纯绿光下，没有可供漫射的红色波段，
衬衫吸收所有波段，呈现黑色。在品红光
$(R, B)$ 下，它只漫射红色波段：仍为红色。白色
物体会漫射它接收到的每一个波段（在绿光下它呈绿色）；
黑色物体在任何光下都吸收所有波段——这也是
黑色衣服在阳光下升温的原因。
\end{example}

\begin{omfigure}
\begin{tikzpicture}[scale=0.9, font=\small]
  \fill[red!80] (0,0) rectangle (1.6,1.0);
  \node at (0.8,0.5) {衬衫};
  \draw[-Stealth, red, thick] (-1.9,1.45) -- (-0.1,0.9);
  \draw[-Stealth, green!65!black, thick] (-1.9,0.95) -- (-0.1,0.62);
  \draw[-Stealth, blue, thick] (-1.9,0.45) -- (-0.1,0.34);
  \draw[-Stealth, red, very thick] (1.7,0.9) -- (3.1,1.45)
    node[above right=-3pt] {红色漫射};
  \node[anchor=west] at (1.7,0.3) {G、B 被吸收};
  \node at (0.8,-0.55) {白光：红色衬衫};
\end{tikzpicture}\qquad
\begin{tikzpicture}[scale=0.9, font=\small]
  \fill[black!85] (0,0) rectangle (1.6,1.0);
  \node[white] at (0.8,0.5) {衬衫};
  \draw[-Stealth, green!65!black, thick] (-1.9,0.95) -- (-0.1,0.62);
  \node[anchor=west] at (1.7,0.75) {无漫射};
  \node[anchor=west] at (1.7,0.3) {G 被吸收};
  \node at (0.8,-0.55) {绿光：黑色衬衫};
\end{tikzpicture}

{\small 同一件红色衬衫在两种照明下：它只能漫射
红色波段——白光中有，绿光中没有。}
\end{omfigure}

\section{光源与色温}

\begin{definition}[白炽光源]\label{def:g11:color-light-sources:incandescent}
因高温而发光的光源——火焰、灯丝、恒星——都以
\emph{白炽}\index{白炽}方式发光。其光谱为
在\cref{ch:g10:light-spectra}中的连续热光谱，其峰值
服从维恩定律（见\cref{prop:g10:light-spectra:wien}），
\[
\lambda_{\max} \times T = \qty{2.90e-3}{m.K},
\]
今后将其作为定量工具：测得峰值，便知温度。
\end{definition}

\begin{example}[效率很低的灯泡]\label{ex:g11:color-light-sources:bulb}
温度 $T = \qty{2700}{K}$ 的钨丝，峰值位于
$\lambda_{\max} = \qty{2.90e-3}{m.K}/\qty{2700}{K} = \qty{1.07e-6}{m}
= \qty{1070}{nm}$，深入红外区：大部分电功率
以不可见辐射的形式散失，仅约 $5\%$ 以光的形式输出——
白炽灯泡实际上主要是一个加热器。
\end{example}

\begin{omfigure}
\begin{tikzpicture}
\begin{axis}[omaxis, width=9.4cm, height=5.4cm,
    xlabel={$\lambda$ (\unit{nm})}, ylabel={光强（任意单位）},
    xmin=0, xmax=2400, ymin=0, ymax=10,
    xtick={400, 800, 1200, 1600, 2000}, ytick=\empty]
  \fill[gray, opacity=0.12] (axis cs:400,0) rectangle (axis cs:800,10);
  \node[gray, anchor=south] at (axis cs:600,8.7) {可见光};
  \addplot[omThm, thick, domain=200:2300, samples=160]
    {3.2e16 * x^(-5) / (exp(2400/x) - 1)};
  \addplot[omDef, thick, domain=250:2300, samples=160]
    {3.84e17 * x^(-5) / (exp(4800/x) - 1)};
  \draw[dashed, omThm] (axis cs:483,0) -- (axis cs:483,8.5);
  \draw[dashed, omDef] (axis cs:967,0) -- (axis cs:967,3.2);
  \node[omThm, anchor=west] at (axis cs:520,9.0) {$\qty{6000}{K}$};
  \node[omDef, anchor=west] at (axis cs:1010,3.4) {$\qty{3000}{K}$ ($\times 12$)};
\end{axis}
\end{tikzpicture}

{\small $\qty{6000}{K}$ 与 $\qty{3000}{K}$ 的热光谱，峰值
$\lambda_{\max}$ 处以虚线标出，分别为 $483$ 与 $\qty{967}{nm}$；较冷的曲线，
放大十二倍，否则将紧贴横轴。}
\end{omfigure}

\begin{definition}[光谱灯与激光器]\label{def:g11:color-light-sources:lamps}
\emph{光谱灯}\index{光谱灯}——一种低压气体，
经放电激发，如钠、汞、氖——发射线状光谱
（见\cref{ch:g10:light-spectra}）：仅有少数孤立波长，其间空无一物。
\emph{激光器}\index{激光器}则更进一步：只发出单一
波长的单色光，且集中在一束狭窄的光束中。
\end{definition}

\begin{definition}[发光二极管]\label{def:g11:color-light-sources:led}
\emph{发光二极管}\index{发光二极管}（LED）发出一个
\emph{准单色的}\index{准单色的}波段，宽度仅几十纳米，
窄到足以呈现某种确定颜色，远比激光谱线宽。白色
LED 实为蓝色 LED（峰值约 $\qty{450}{nm}$）外覆
\emph{荧光粉}\index{荧光粉}，一种吸收部分蓝光的材料，
并以较宽的黄色波段重新发射：蓝色与黄色
互补（见\cref{prop:g11:color-light-sources:complementary}），
两者混合后呈现白色。
\end{definition}

\begin{definition}[色温]\label{def:g11:color-light-sources:colortemp}
类白色光源的\emph{色温}\index{色温}，是指与其光色
色调相符的白炽体的温度，常见值为：
蜡烛 $\qty{1800}{K}$、白炽灯泡 $\qty{2700}{K}$、
卤素灯 $\qty{3400}{K}$、正午阳光 $\qty{5500}{K}$、阴天日光
$\qty{6500}{K}$。它是一个颜色标签，而非温度计的读数——一个
$\qty{6500}{K}$ 的 LED 仍处于室温——而“冷”“暖”的用法
却恰好相反：低色温呈“暖”的橙色，高色温呈“冷”的
偏蓝色。
\end{definition}

\begin{omfigure}
\begin{tikzpicture}[font=\small]
  \shade[left color=orange!95!red, right color=white] (0,0) rectangle (5.4,0.7);
  \shade[left color=white, right color=blue!25] (5.4,0) rectangle (9,0.7);
  \draw (0,0) rectangle (9,0.7);
  \draw[-Stealth] (0,-0.5) -- (9.6,-0.5)
    node[right] {色温 (\unit{K})};
  \foreach \x/\l in {0.5/1800, 2.0/2700, 3.1/3400, 6.5/5500, 8.2/6500}
    \draw (\x,-0.4) -- (\x,-0.6) node[below] {\l};
  \foreach \x/\n in {0.5/烛光，2.0/白炽灯，3.1/卤素灯，6.5/阳光，
      8.2/阴天}
    \node[above] at (\x,0.7) {\n};
\end{tikzpicture}

{\small 常见光源的色温标尺：“暖”色调的橙色
位于低 $T$ 处，“冷”的偏蓝色调位于高 $T$ 处。}
\end{omfigure}

\section{练习}

\begin{exercise}[$\star$]\label{exo:g11:color-light-sources:1}
求白色屏幕同时接收下列光时所感知的颜色：
(a) 红光和绿光；(b) 绿光和蓝光；(c) 三种
原色光。哪一种光与蓝光互补？
\end{exercise}

\begin{exercise}[$\star$]\label{exo:g11:color-light-sources:2}
屏幕需点亮哪些子像素（R、G、B）才能显示以下颜色：黄色、
品红色、白色、黑色、橙色？
\end{exercise}

\begin{exercise}[$\star$]\label{exo:g11:color-light-sources:3}
白光通过一块品红色滤光片。哪一波段被吸收，
透出的是什么颜色？再在其后加一块青色滤光片，此时透出的又是什么颜色？
\end{exercise}

\begin{exercise}[$\star$]\label{exo:g11:color-light-sources:4}
在白光下，一面旗帜上有一条黄色条纹。它漫射哪些
波段，吸收哪些波段？它在 (a) 红光、(b) 蓝光下分别呈现什么颜色？
\end{exercise}

\begin{exercise}[$\star$]\label{exo:g11:color-light-sources:5}
某恒星的连续光谱在 $\lambda_{\max} = \qty{580}{nm}$ 处达到峰值。
求其表面温度。
\end{exercise}

\begin{exercise}[$\star\star$]\label{exo:g11:color-light-sources:6}
蜡烛火焰的温度约为 $\qty{1800}{K}$。(a) 求
$\lambda_{\max}$，并指出其所属的光谱区域。(b) 火焰看起来却
呈橙色，试加以解释。
\end{exercise}

\begin{exercise}[$\star\star$]\label{exo:g11:color-light-sources:7}
一名骑车人身穿红色夹克，手戴青色手套。
利用\cref{met:g11:color-light-sources:bookkeeping}，求出
每件物品在 (a) 白光、(b) 红光、(c) 绿光下的表观颜色。
\end{exercise}

\begin{exercise}[$\star\star$]\label{exo:g11:color-light-sources:8}
一条老街由低压钠灯照明（钠灯的光谱仅有一条
$\qty{589}{nm}$ 处的谱线）。(a) 这是哪一类光谱？在其照射下，
\emph{任何}物体能呈现哪些颜色？(b) 一辆汽车在这些
灯下呈黑色，而在日光下呈蓝色：试加以解释。
\end{exercise}

\begin{exercise}[$\star\star$]\label{exo:g11:color-light-sources:9}
一支绿色激光笔发出 $\qty{532}{nm}$ 的光。(a) 其光束穿过
一个棱镜：将结果与白光光束的结果相比较，并说明所体现的
性质。(b) 随后它遇到一块红色滤光片：会透出什么？
\end{exercise}

\begin{exercise}[$\star\star$]\label{exo:g11:color-light-sources:10}
描述白色 LED（蓝色发光二极管加荧光粉）的光谱，并
解释为何人眼感知其光为白色。为何深红色物体在廉价白色 LED
下显得比在日光下更暗淡？
\end{exercise}

\begin{exercise}[$\star\star$]\label{exo:g11:color-light-sources:11}
求白炽体的 $\lambda_{\max}$，分别取 (a) 蜡烛的
色温（$\qty{1800}{K}$）与 (b) 阴天日光的色温
（$\qty{6500}{K}$）时的值。哪一种光被称为“暖”光？请加以评论。
\end{exercise}

\begin{exercise}[$\star\star\star$]\label{exo:g11:color-light-sources:12}
一盏钠灯（$\qty{589}{nm}$）与一块显示黄色的屏幕
（子像素发出 $\qty{610}{nm}$ 与 $\qty{540}{nm}$ 的光）看上去
颜色完全相同。(a) 这样的两种光分别称为什么，它们的存在
证明了关于颜色的什么？(b) 提出两种不同的实验以区分它们，
并预测每种实验的结果。
\end{exercise}

\begin{exercise}[$\star\star\star$]\label{exo:g11:color-light-sources:13}
一台剧场卤素聚光灯的色温为 $\qty{3200}{K}$。(a) 求
$\lambda_{\max}$，判断该光比日光更暖还是更冷？(b) 为模拟日光，
在聚光灯前覆一层淡蓝色滤色片。滤色片滤去了什么，
舞台为何变暗？(c) 解释为何没有任何滤色片
能够反而\emph{增强}聚光灯光谱的蓝端。
\end{exercise}

\begin{exercise}[$\star\star\star$]\label{exo:g11:color-light-sources:14}
一名屠夫用色温为
$\qty{2700}{K}$ 的发光二极管照亮肉类陈列柜，隔壁的鱼贩则用 $\qty{6500}{K}$ 的灯。
(a) 描述每种照明的色调，及其对红肉与
白鱼的影响。(b) 两位商贩中是否有人改变了所售商品？请指出
变化究竟发生在何处（见\cref{def:g11:color-light-sources:objectcolor}）。
\end{exercise}

\begin{exercise}[$\star\star\star$]\label{exo:g11:color-light-sources:15}
某镇紧邻一座天文台，须重新点亮街道，
在低压钠灯与宽谱白光发光二极管之间犹豫不决。
(a) 天文学家更倾向于哪种方案？为什么？（想一想，一片
滤光片能从他们的图像中滤去什么。）(b) 司机与行人
更倾向于哪种方案？请用物体的颜色说明理由。
(c) 提出一个折中方案，并用色温加以表述。
\end{exercise}

\section{问题：颜色的工程}

\begin{problem}[{周末问题——舞台灯光、钠灯街道与
白光发光二极管：颜色如何被设计出来，从可调出任意色调的三盏聚光灯，
到点亮本页的那片芯片}]
\label{pb:g11:color-light-sources:1}
一位灯光设计师只拥有红、绿、蓝三色聚光灯；一名街道工程师
拥有一盏只发出单一波长光的灯；一位芯片制造商
拥有一只仅发蓝光的二极管。三人都在出售颜色。本题依次研究他们的行当——
舞台上的加色混合、街道上的单色经济，
芯片中的荧光粉“炼金术”，最后
为卧室和珠宝店的橱窗选择光源。

\medskip
\noindent\textbf{第一部分——舞台：用三盏聚光灯“作画”。}
\begin{enumerate}
  \item 三盏聚光灯在白色幕布上重叠。说出两两重叠区域
        的颜色及三者重叠区域的颜色。
  \item 设计师如何调出橙色？又如何调出淡粉色？（定性描述
        各聚光灯的功率。）
  \item 一名舞者站在幕布前，仅由红色和绿色
        聚光灯照亮。幕布上出现两个有色的影子，各是什么颜色？请解释。
  \item 幕布重新涂成深蓝色。它在仅有红色聚光灯照射下呈现何种颜色？
        三盏聚光灯全部打开时又是什么颜色？
  \item 为什么三盏聚光灯能重现“几乎所有”颜色？人体的哪一
        部位决定了数目为三？
\end{enumerate}

\medskip
\noindent\textbf{第二部分——服装：减色的作用。}
\begin{enumerate}[resume]
  \item 一件服装在白光下呈黄色。其染料吸收哪一个
        波段的光？
  \item 预测它在红色聚光灯、绿色聚光灯和蓝色聚光灯
        下分别呈现的颜色。
  \item 设计师希望只需轻拨一下开关，服装便能由亮
        变黑。拨动这一下开关，应切换到哪一盏聚光灯？
  \item 一片品红色滤色片被推到照亮
        黄色服装的白色聚光灯前。追踪各波段
        （见\cref{met:g11:color-light-sources:bookkeeping}）：服装呈何种颜色？
  \item 舞台上的滤色片为何做成青色、品红和黄色，
        而不是红、绿、蓝？（红色滤色片会对白色
        聚光灯的功率产生什么影响？）
\end{enumerate}

\medskip
\noindent\textbf{第三部分——钠灯街道：单色经济。}
\begin{enumerate}[resume]
  \item 一盏低压钠灯几乎只在
        $\qty{589}{nm}$ 处发出一条谱线。这是哪一类光谱？整条街道在它照射下
        呈现怎样的样子？
  \item 在此灯照射下，写出白墙、
        黄色路面标线、红色汽车和蓝色汽车的表观颜色。
  \item 天文学家为何喜爱这些灯？（一片滤光片能对
        全城的灯光辉光产生什么作用？）
  \item 有人建议改用色温为 $\qty{2700}{K}$ 的钨丝街灯灯泡。
        求其 $\lambda_{\max}$，并说明其大部分功率分布在何处？
  \item 钠灯约将其电功率的 $30\%$
        转化为可见光，钨丝灯泡约为 $5\%$。
        每输入 $\qty{100}{W}$ 电功率，两种灯各能输出多少瓦特的光？
\end{enumerate}

\medskip
\noindent\textbf{第四部分——白光发光二极管与光源的选择。}
\begin{enumerate}[resume]
  \item 白光发光二极管是覆有黄色
        荧光粉的蓝色二极管（$\qty{450}{nm}$）。借助互补色解释，
        为何其输出光看起来是白色。
  \item 描述其光谱，并将它与日光的
        连续光谱加以比较。哪一波段明显偏弱？
  \item 推断这种发光二极管对哪些物体的显色效果较差，并将此与
        \cref{exo:g11:color-light-sources:14}中的屠夫联系起来。
  \item 某款发光二极管灯泡标注为“$\qty{2700}{K}$”。其内部是否有任何部件
        的温度为 $\qty{2700}{K}$？标签承诺了什么？求出它所模拟的
        白炽体的 $\lambda_{\max}$。
  \item 结语——各用一句话说明理由，为卧室灯与
        珠宝店橱窗各选择一个色温；然后用一句话
        概括本章的颜色定律。
\end{enumerate}
\end{problem}
